\documentclass[10pt,a4paper]{amsart}
\usepackage{amsmath,amssymb,amsthm,graphicx}
% Retained preamble for the Stacks Project chapter divisors.tex (Divisors).
% The chapter file pinned-original-source.tex carries only \input{preamble} on its first line, so no
% preamble text lives inside the chapter itself.  The Stacks Project uses one single shared preamble for
% every chapter; the snapshot stages it verbatim in sources/stacks-chow/preamble.tex.  The block below is
% copied mechanically, line for line, from that file's lines 14-18 (the Xy-pic setup) and lines 154-217
% (the native theorem declarations together with the mathematical macros).  Nothing was added, removed or
% redefined.  The wrapper class and the amsmath/amssymb/amsthm/graphicx base packages are supplied by the
% assembler, so the class/\externaldocument part of the shared preamble is not repeated here.
% For commutative diagrams we use Xy-pic
\usepackage[all]{xy}

% We use 2cell for 2-commutative diagrams.
\xyoption{2cell}
% Theorem environments.
%
\theoremstyle{plain}
\newtheorem{theorem}[subsection]{Theorem}
\newtheorem{proposition}[subsection]{Proposition}
\newtheorem{lemma}[subsection]{Lemma}

\theoremstyle{definition}
\newtheorem{definition}[subsection]{Definition}
\newtheorem{example}[subsection]{Example}
\newtheorem{exercise}[subsection]{Exercise}
\newtheorem{situation}[subsection]{Situation}

\theoremstyle{remark}
\newtheorem{remark}[subsection]{Remark}
\newtheorem{remarks}[subsection]{Remarks}

\numberwithin{equation}{subsection}

% Macros
%
\def\lim{\mathop{\mathrm{lim}}\nolimits}
\def\colim{\mathop{\mathrm{colim}}\nolimits}
\def\Spec{\mathop{\mathrm{Spec}}}
\def\Hom{\mathop{\mathrm{Hom}}\nolimits}
\def\Ext{\mathop{\mathrm{Ext}}\nolimits}
\def\SheafHom{\mathop{\mathcal{H}\!\mathit{om}}\nolimits}
\def\SheafExt{\mathop{\mathcal{E}\!\mathit{xt}}\nolimits}
\def\Sch{\mathit{Sch}}
\def\Mor{\mathop{\mathrm{Mor}}\nolimits}
\def\Ob{\mathop{\mathrm{Ob}}\nolimits}
\def\Sh{\mathop{\mathit{Sh}}\nolimits}
\def\NL{\mathop{N\!L}\nolimits}
\def\CH{\mathop{\mathrm{CH}}\nolimits}
\def\proetale{{pro\text{-}\acute{e}tale}}
\def\etale{{\acute{e}tale}}
\def\QCoh{\mathit{QCoh}}
\def\Ker{\mathop{\mathrm{Ker}}}
\def\Im{\mathop{\mathrm{Im}}}
\def\Coker{\mathop{\mathrm{Coker}}}
\def\Coim{\mathop{\mathrm{Coim}}}

% Boxtimes
%
\DeclareMathSymbol{\boxtimes}{\mathbin}{AMSa}{"02}

%
% Macros for moduli stacks/spaces
%
\def\QCohstack{\mathcal{QC}\!\mathit{oh}}
\def\Cohstack{\mathcal{C}\!\mathit{oh}}
\def\Spacesstack{\mathcal{S}\!\mathit{paces}}
\def\Quotfunctor{\mathrm{Quot}}
\def\Hilbfunctor{\mathrm{Hilb}}
\def\Curvesstack{\mathcal{C}\!\mathit{urves}}
\def\Polarizedstack{\mathcal{P}\!\mathit{olarized}}
\def\Complexesstack{\mathcal{C}\!\mathit{omplexes}}
% \Pic is the operator that assigns to X its picard group, usage \Pic(X)
% \Picardstack_{X/B} denotes the Picard stack of X over B
% \Picardfunctor_{X/B} denotes the Picard functor of X over B
\def\Pic{\mathop{\mathrm{Pic}}\nolimits}
\def\Picardstack{\mathcal{P}\!\mathit{ic}}
\def\Picardfunctor{\mathrm{Pic}}
\def\Deformationcategory{\mathcal{D}\!\mathit{ef}}

% --- appended 2026-10-12 (Chief): the project own definitions of the reference/slogan/history
% display environments, copied mechanically from the shared Stacks preamble lines 8-12
% (\usepackage{verbatim} plus the three \newenvironment lines, which the chapter preamble keeps
% verbatim from the project preamble). Nothing added, removed or redefined: these environments
% carry no mathematical content and the project itself typesets them as comments. Without these
% three lines a counted statement whose own text uses them fails with Environment slogan undefined.
\usepackage{verbatim}
\newenvironment{reference}{\comment}{\endcomment}
\newenvironment{slogan}{\comment}{\endcomment}
\newenvironment{history}{\comment}{\endcomment}
\usepackage{cleveref}
\crefname{originaldefinition}{original definition}{original definitions}
\Crefname{originaldefinition}{Original definition}{Original definitions}
\makeatletter
\newlabel{definition-admissible-blowup}{{080K}{source}{}{sourcepointer-1}{}}
\newlabel{definition-admissible-blowup@cref}{{[originaldefinition][0][]080K}{[1][0][]source}{}{}{}}
\makeatother
\crefname{originallemma}{original lemma}{original lemmas}
\Crefname{originallemma}{Original lemma}{Original lemmas}
\makeatletter
\newlabel{lemma-fitting-ideal-finite-locally-free}{{0C3G}{source}{}{sourcepointer-2}{}}
\newlabel{lemma-fitting-ideal-finite-locally-free@cref}{{[originallemma][0][]0C3G}{[1][0][]source}{}{}{}}
\makeatother
\crefname{originallemma}{original lemma}{original lemmas}
\Crefname{originallemma}{Original lemma}{Original lemmas}
\makeatletter
\newlabel{lemma-fitting-ideal-of-finitely-presented}{{0C3E}{source}{}{sourcepointer-3}{}}
\newlabel{lemma-fitting-ideal-of-finitely-presented@cref}{{[originallemma][0][]0C3E}{[1][0][]source}{}{}{}}
\makeatother
\crefname{originallemma}{original lemma}{original lemmas}
\Crefname{originallemma}{Original lemma}{Original lemmas}
\makeatletter
\newlabel{lemma-strict-transform-blowup-fitting-ideal-locally-free}{{0CZQ}{source}{}{sourcepointer-4}{}}
\newlabel{lemma-strict-transform-blowup-fitting-ideal-locally-free@cref}{{[originallemma][0][]0CZQ}{[1][0][]source}{}{}{}}
\makeatother
\crefname{originallemma}{original lemma}{original lemmas}
\Crefname{originallemma}{Original lemma}{Original lemmas}
\makeatletter
\newlabel{modules-lemma-pullback-finite-presentation}{{01BQ}{source}{}{sourcepointer-5}{}}
\newlabel{modules-lemma-pullback-finite-presentation@cref}{{[originallemma][0][]01BQ}{[1][0][]source}{}{}{}}
\makeatother
\crefname{originallemma}{original lemma}{original lemmas}
\Crefname{originallemma}{Original lemma}{Original lemmas}
\makeatletter
\newlabel{more-algebra-lemma-fitting-ideals-and-pd1}{{0F7M}{source}{}{sourcepointer-6}{}}
\newlabel{more-algebra-lemma-fitting-ideals-and-pd1@cref}{{[originallemma][0][]0F7M}{[1][0][]source}{}{}{}}
\makeatother
\title{Blowing up the r-th Fitting ideal makes a module finite locally free of rank r on a scheme theoretically dense open}
\author{The Stacks Project Authors}
\date{stacks-project:divisors.tex:local-snapshot-20261012; GNU FDL 1.2 or later}
\begin{document}
\maketitle
\noindent\textit{Editorial note.} All delivered mathematical text is the author's own, copied from the pinned Stacks Project chapter file pinned-original-source.tex; no statement or proof step was rewritten, reordered or removed.\par\medskip
\noindent\textit{Editorial note.} Following the multi-claim convention for this source, source\_url is the source-level URL https://stacks.math.columbia.edu/ and the source label lemma-blowup-fitting-ideal is carried in claim\_key and here rather than in a per-claim tag URL.\par\medskip
\noindent\textit{Editorial note.} The preamble is the chapter's shared Stacks Project preamble, retained as sources/stacks-divisors/retained-preamble.tex: a line-for-line excerpt of the shared Stacks preamble staged with this snapshot. No macro was invented or redefined, and this card adds no preamble text of its own.\par\medskip
\noindent\textit{Editorial note.} Statement span 9339-9356 and proof span 9358-9397 come from the halfA pool card operations/pool-cards/stacks-divisors-halfA.json and were re-read in the pinned chapter file before this card was written.\par\medskip
\noindent\textit{Editorial note.} The proof contains the author's parenthesis noting that the complement of the Fitting subscheme may be bigger than the original open, which is why the density step is needed; it is preserved verbatim.\par\medskip
\noindent\textit{Editorial note.} 6 cross-references occur inside the counted spans; they are listed in external\_references as navigation pointers only.\par\medskip
\noindent\textit{Allowed standard prerequisites (editorial).}\begin{itemize}
\item The r-th Fitting ideal of a finitely presented module is of finite type, and its restriction to a locus where the module is finite locally free of rank r is the whole structure sheaf, while the (r-1)-st Fitting ideal restricts to zero there. Cited at proof lines 9358-9364 and 9369-9372. Gives admissibility of the blowing up and shows that the smaller Fitting ideal vanishes identically because the open is scheme theoretically dense.
\item The definition of an admissible blowing up, and the statement that a module becomes finite locally free of the expected rank after blowing up the appropriate Fitting ideal. Cited at proof lines 9364-9366 and 9378-9382. Supplies the notion of admissibility and the conclusion about the strict transform.
\item The pullback of a finitely presented module is finitely presented, cited from the modules chapter. Cited at proof lines 9390-9392. Shows that the kernel of the surjection onto the strict transform is finitely presented once the pullback is split.
\item The vanishing of the higher Tor of a module whose Fitting ideals behave as here, cited from the more-algebra chapter. Cited at proof lines 9395-9397. Gives the statement about perfectness and tor dimension.
\end{itemize}\medskip
\ifdefined\hypertarget\hypertarget{sourcepointer-1}{}\fi
\noindent\textit{Source reference.} Original definition 080K: focused-seeds/AGP/sources/stacks-snapshot-20261011/tags.txt, source lines 10185--10185. Same-chapter reference used in the retained author text at source line(s) 9364. It is supplied only as a navigation pointer so that the retained text compiles; the counted item delivers the author's own statement and the author's complete proof, and no part of the referenced item is copied in or substituted for the counted argument.\par\medskip
\ifdefined\hypertarget\hypertarget{sourcepointer-2}{}\fi
\noindent\textit{Source reference.} Original lemma 0C3G: focused-seeds/AGP/sources/stacks-snapshot-20261011/tags.txt, source lines 15191--15191. Same-chapter reference used in the retained author text at source line(s) 1441, 9362, 9367, 9372. It is supplied only as a navigation pointer so that the retained text compiles; the counted item delivers the author's own statement and the author's complete proof, and no part of the referenced item is copied in or substituted for the counted argument.\par\medskip
\ifdefined\hypertarget\hypertarget{sourcepointer-3}{}\fi
\noindent\textit{Source reference.} Original lemma 0C3E: focused-seeds/AGP/sources/stacks-snapshot-20261011/tags.txt, source lines 15189--15189. Same-chapter reference used in the retained author text at source line(s) 1450, 9360. It is supplied only as a navigation pointer so that the retained text compiles; the counted item delivers the author's own statement and the author's complete proof, and no part of the referenced item is copied in or substituted for the counted argument.\par\medskip
\ifdefined\hypertarget\hypertarget{sourcepointer-4}{}\fi
\noindent\textit{Source reference.} Original lemma 0CZQ: focused-seeds/AGP/sources/stacks-snapshot-20261011/tags.txt, source lines 16285--16285. Same-chapter reference used in the retained author text at source line(s) 9378. It is supplied only as a navigation pointer so that the retained text compiles; the counted item delivers the author's own statement and the author's complete proof, and no part of the referenced item is copied in or substituted for the counted argument.\par\medskip
\ifdefined\hypertarget\hypertarget{sourcepointer-5}{}\fi
\noindent\textit{Source reference.} Original lemma 01BQ: focused-seeds/AGP/sources/stacks-snapshot-20261011/tags.txt, source lines 1735--1735. Same-chapter reference used in the retained author text at source line(s) 9391. It is supplied only as a navigation pointer so that the retained text compiles; the counted item delivers the author's own statement and the author's complete proof, and no part of the referenced item is copied in or substituted for the counted argument.\par\medskip
\ifdefined\hypertarget\hypertarget{sourcepointer-6}{}\fi
\noindent\textit{Source reference.} Original lemma 0F7M: focused-seeds/AGP/sources/stacks-snapshot-20261011/tags.txt, source lines 19012--19012. Same-chapter reference used in the retained author text at source line(s) 9396. It is supplied only as a navigation pointer so that the retained text compiles; the counted item delivers the author's own statement and the author's complete proof, and no part of the referenced item is copied in or substituted for the counted argument.\par\medskip
\begin{lemma}
\label{lemma-blowup-fitting-ideal}
Let $X$ be a scheme. Let $\mathcal{F}$ be a finitely presented
$\mathcal{O}_X$-module. Let $U \subset X$ be a scheme theoretically
dense open such that $\mathcal{F}|_U$ is finite locally free of
constant rank $r$. Then
\begin{enumerate}
\item the blowup $b : X' \to X$ of $X$ in the $r$th Fitting
ideal of $\mathcal{F}$ is $U$-admissible,
\item the strict transform $\mathcal{F}'$ of $\mathcal{F}$
with respect to $b$ is locally free of rank $r$,
\item the kernel $\mathcal{K}$ of the surjection
$b^*\mathcal{F} \to \mathcal{F}'$ is
finitely presented and $\mathcal{K}|_U = 0$,
\item $b^*\mathcal{F}$ and $\mathcal{K}$ are perfect
$\mathcal{O}_{X'}$-modules of tor dimension $\leq 1$.
\end{enumerate}
\end{lemma}

\begin{proof}
The ideal $\text{Fit}_r(\mathcal{F})$ is of finite type
by Lemma \ref{lemma-fitting-ideal-of-finitely-presented}
and its restriction to $U$ is equal to $\mathcal{O}_U$ by
Lemma \ref{lemma-fitting-ideal-finite-locally-free}.
Hence $b : X' \to X$ is $U$-admissible, see
Definition \ref{definition-admissible-blowup}.

\medskip\noindent
By Lemma \ref{lemma-fitting-ideal-finite-locally-free}
the restriction of $\text{Fit}_{r - 1}(\mathcal{F})$
to $U$ is zero, and since $U$ is scheme theoretically dense
we conclude that $\text{Fit}_{r - 1}(\mathcal{F}) = 0$
on all of $X$. Thus it follows from
Lemma \ref{lemma-fitting-ideal-finite-locally-free}
that $\mathcal{F}$ is locally free of rank $r$
on the complement of subscheme cut out by the $r$th
Fitting ideal of $\mathcal{F}$ (this complement may
be bigger than $U$ which is why we had to do this step
in the argument). Hence by
Lemma \ref{lemma-strict-transform-blowup-fitting-ideal-locally-free}
the strict transform
$$
b^*\mathcal{F} \longrightarrow \mathcal{F}'
$$
is locally free of rank $r$. The kernel $\mathcal{K}$
of this map is supported on the exceptional divisor
of the blowup $b$ and hence $\mathcal{K}|_U = 0$.
Finally, since $\mathcal{F}'$ is finite locally free
and since the displayed arrow is surjective, we can
locally on $X'$ write $b^*\mathcal{F}$ as the
direct sum of $\mathcal{K}$ and $\mathcal{F}'$.
Since $b^*\mathcal{F}'$ is finitely presented
(Modules, Lemma \ref{modules-lemma-pullback-finite-presentation})
the same is true for $\mathcal{K}$.

\medskip\noindent
The statement on tor dimension follows from
More on Algebra, Lemma \ref{more-algebra-lemma-fitting-ideals-and-pd1}.
\end{proof}

\end{document}
